\begin{abstract}
Proactive emotional care requires more than recognizing whether a user looks negative. A practical agent must estimate latent affective state, decide whether the current moment is receptive to intervention, and choose a support strategy that is helpful without becoming intrusive. We study this problem as a staged pipeline with an upstream tri-modal emotion encoder and a downstream timing-and-strategy module, and implement a reproducible prototype benchmark around synthetic data and weak labels from deployed logs.

Our current release uses three offline data sources plus a pending real-pilot protocol: a timing synthetic split with 24{,}000/2{,}400/2{,}400 windows, a tri-modal synthetic emotion split with 120{,}000/12{,}000/12{,}000 windows, and a weak-label split with 286/61/62 windows extracted from 656 historical emotion events and 2{,}838 chat messages. In formal dual-A800 experiments, the tri-modal face+motion+audio encoder reaches 97.59\% accuracy and 97.58\% emotion Macro-F1 on the held-out synthetic fine-emotion benchmark, while single-modality ablations collapse to 24.33\%, 8.33\%, and 31.11\% accuracy. Downstream, the exploratory joint temporal timing model reaches 86.90\% balanced accuracy and 81.06\% timing Macro-F1 on the held-out synthetic timing split. On the held-out 62-window weaklabel split, however, the same joint model reaches only 36.29\% timing Macro-F1 and 18.06\% strategy Macro-F1, exposing a substantial synthetic-to-log domain gap.

These results support a narrower claim than deployment readiness: the repository now provides a complete cold-start pipeline, a usable weak-label bootstrapping route, and a transparent set of baselines for the proactive care timing problem. It does not yet establish raw-video end-to-end performance, real-user benefit, or strong log-domain robustness.
\end{abstract}
